The answer is difficult to obtain because BV presents two coupled ensembles. Its four nitrogens and two lactam carbonyl groups can accommodate protons in several arrangements, including keto--enol tautomers, and each arrangement can adopt multiple conformations.\cite{villegas-escobarKetoEnolTautomerism2021a,tongManyRoadsSeam2026} Rotations about the three methine bridges produce distinct conformers, including the locked (\emph{ZZZsss}) and open (\emph{ZZZssa}) helices, while fluctuations within either basin generate structures with different helicities (\textcolor{red}{Figure X: BV isomers and helicity, WIP}). The free energy of removing a proton from one site therefore depends on both the locations of the remaining protons and the conformation of the molecule. A measured \pka{} averages over these coupled protonation and conformational ensembles. Even rare conformers can contribute strongly if their acidity compensates for their low population; in carbonic acid, for example, a conformer representing only about 12\% of the ensemble releases most of the protons.\cite{zhaoInitioDeepNeural2025} Most practical computational approaches avoid calculating the absolute transfer of a proton from a specific site into bulk water and instead use relative binding free energies, empirical reference compounds, thermodynamic cycles, or continuum electrostatics.\cite{chengRedoxPotentialsAcidity2014,navoComputerPredictionKa2021,pliegoThermodynamicCyclesCalculation2003,casasnovasTheoreticalKaCalculations2014} These methods can rank protonation states or predict relative shifts, but they do not directly describe site-resolved proton release, competing proton endpoints, or its coupling to the conformational ensemble. Calculations on bilins have consequently focused on relative proton affinities in continuum or protein-like environments,\cite{borgRelativeGroundExcitedState2007,corbellaDeterminationProtonationPreferences2018} and, to our knowledge, no absolute site-resolved \pka{} calculation of free BV in explicit water has been reported.

In this work, we calculate this equilibrium using reactive divide-and-conquer density functional tight binding (DC-DFTB) metadynamics in explicit water.\cite{barducciMetadynamics2011,nishizawaThreePillarsAchieving2016,nakaiDivideandConquerTypeDensityFunctionalTightBinding2016,nishimuraHierarchicalParallelizationDivideandconquer2020} DFTB-based metadynamics reproduced the measured \pka{} values of 34 small aqueous amines with a mean absolute deviation of 0.09 units,\cite{Sakti2018} while the divide-and-conquer formulation keeps the computational cost close to linear in system size, making a solvated BV system of approximately one thousand atoms accessible.\cite{nishizawaThreePillarsAchieving2016,nakaiDivideandConquerTypeDensityFunctionalTightBinding2016} In an absolute deprotonation calculation, water is not merely the environment but a reactant that accepts and transports the released proton. The excess-proton defect moves through reorganizing hydrogen-bonded water wires by Grotthuss structural diffusion, rather than as one identifiable hydrogen atom.\cite{agmonGrotthussMechanism1995,agmonProtonsHydroxideIons2016,cukiermanTuGrotthussOther2006,hassanaliProtonTransferWater2013,miyakeGrotthussMechanismsProton2016,popovSearchGrotthussMechanism2023} We promote this process with a simple coordination-number collective variable, which has been shown to distinguish the acid and base states when the resulting trajectories are analyzed correctly.\cite{Tummanapelli2015Polyprotic,Tummanapelli2015Amino,Sakti2018} Because the variable does not prescribe the proton acceptor or endpoint, we subsequently follow the excess-charge defect through the solvent and retain only trajectories in which it reaches bulk water, separating completed deprotonation from return, intramolecular transfer, interfacial trapping, and counterion recombination.

Complete bulk-reaching deprotonation trajectories are rare for a system of this size, so we use one directly calculated absolute \pka{} as the thermodynamic anchor and obtain the remaining acidities from the relative free energies connecting the BV microstates. DC-DFTB umbrella sampling determines the free energies of proton transfer between adjacent rings and of the ring-D rotation that interconverts the locked and open helices. Thermodynamic cycles then combine these relative free energies with the absolute anchor to determine the tautomer populations and macroscopic \pka{} of both conformational ensembles. Three smaller molecules test the individual components of this strategy: histidine, with two competing sites of known microscopic \pka{};\cite{kyte_structure_2006,tanokura1HNMRStudyTautomerism1983} propionic acid, whose \emph{syn} and \emph{anti} conformers close a cycle between deprotonation and rotation;\cite{macoas_internal_2005} and malonaldehyde, a symmetric reference for intramolecular proton transfer.\cite{baughcumMicrowaveSpectroscopicStudy1981,wangFulldimensionalQuantumCalculations2008}
